\section{Introduction}
% =====================================================================

Les relevés de galaxies comme DESI cartographient la structure à grande échelle de l'Univers sur des volumes assez grands pour que les erreurs statistiques ne soient plus la principale limitation. La façon dont les galaxies sont distribuées nous renseigne sur la croissance des structures, sur la quantité de matière, et sur le rôle de l'énergie noire et des neutrinos. La question est de savoir quelle part de cette information nos statistiques extraient réellement.

La plupart des analyses compriment le champ de galaxies en statistiques à deux points, la fonction de corrélation ou le spectre de puissance. Pour un champ gaussien, cela ne perd rien, et aux grandes échelles cela fonctionne très bien. Mais la gravité rend l'Univers tardif fortement non gaussien : la matière s'effondre en halos, se rassemble en filaments et en nœuds, et vide les vides, et les vitesses particulières déforment l'image dans l'espace des redshifts. La fonction à deux points ne voit pas tout cela. Les statistiques construites sur la toile cosmique récupèrent une partie de l'information manquante : \cite{Sunseri} a montré avec Quijote que découper le champ de densité en nœuds, filaments, feuillets et vides, puis combiner les spectres de puissance de ces environnements, resserre nettement les contraintes, surtout sur la masse des neutrinos.

Pour les galaxies, deux familles de statistiques d'environnement ont été développées dans ce but. Le clustering par découpage en densité (density split) \cite{Paillas2023, Paillas2024} classe des points aléatoires selon la densité de galaxies lissée autour d'eux, les répartit en quantiles, et mesure l'auto-corrélation de chaque quantile ainsi que sa cross-corrélation avec les galaxies ; combiné à un émulateur issu de simulations \cite{CuestaLazaro2024}, il améliore de plusieurs fois les contraintes $\Lambda$CDM par rapport à la fonction à deux points. Les statistiques de vides, combinées au clustering des galaxies dans un modèle direct de la connexion galaxie--halo, lèvent aussi des dégénérescences entre cosmologie et formation des galaxies \cite{Salcedo}.

ASTRA \cite{ASTRA} emprunte une autre voie. Il triangule les galaxies avec un catalogue aléatoire uniforme et attribue à chaque point un environnement à partir de la fraction de galaxies parmi ses voisins de Delaunay : il n'y a pas d'échelle de lissage, et l'environnement s'adapte à la densité locale. ASTRA peut classer les points en vides, feuillets, filaments et nœuds, ou les ordonner de façon continue.

\paragraph{Objectif.} Notre but est de quantifier de combien les statistiques d'environnement ASTRA améliorent les contraintes sur les paramètres cosmologiques ($\Omega_m$ via $\omega_{\rm cdm}$, $\sigma_8$, les paramètres d'énergie noire $w_0$ et $w_a$, et idéalement la masse des neutrinos) par rapport à la fonction à deux points standard des galaxies, et comment elles se comparent au density split lorsque les deux sont mesurés sur les mêmes mocks, avec les mêmes points aléatoires et les mêmes estimateurs. Pour cela, nous construisons des jeux de données dédiés de covariances et de dérivées à partir de suites de catalogues simulés, qui donnent des prévisions de Fisher des incertitudes sur les paramètres, et nous entraînons un émulateur par réseau de neurones des statistiques, que nous utilisons pour un test d'inférence de bout en bout. La masse des neutrinos s'est révélée hors de portée des simulations disponibles (Section~\ref{sec:quijote}) ; tous les autres paramètres sont couverts.

\paragraph{Données.} Nous utilisons trois suites de simulations, chacune pour ce que les autres ne peuvent pas fournir (Section~\ref{sec:sims}) : la suite d'émulation AbacusSummit de mocks de galaxies de type DESI sur 85 cosmologies, sur laquelle sont construits les mesures et l'émulateur ; les mocks officiels de DESI, dont les 1786 boîtes indépendantes donnent une covariance fiable ; et les simulations Quijote, dont les simulations appariées donnent des dérivées propres. Obtenir des dérivées et des covariances fiables a été la principale difficulté du travail, et nous décrivons comment nos choix ont été faits et comment les résultats ont évolué.

\paragraph{Plan.} La Section~\ref{sec:env} présente les deux estimateurs d'environnement, la Section~\ref{sec:clustering} les mesures de clustering, la Section~\ref{sec:framework} la covariance, les dérivées et le formalisme de Fisher, la Section~\ref{sec:choices} les choix d'analyse successifs, la Section~\ref{sec:results} les résultats (mocks de type DESI, Quijote, classes ASTRA, inférence), la Section~\ref{sec:robustness} les tests de robustesse, et la Section~\ref{sec:discussion} la discussion. Les annexes expliquent comment reproduire l'article (\ref{app:repro}), donnent les détails mathématiques (\ref{app:maths}), et retracent le cheminement méthodologique (\ref{app:path}).
